\documentclass[10pt,a4paper]{amsart}
\usepackage[margin=19mm]{geometry}
\usepackage{amsmath,amssymb,mathtools}
\usepackage[scr=rsfs,cal=euler]{mathalfa}
\usepackage{xfrac}
\usepackage[unicode,hidelinks,bookmarks=false]{hyperref}
\newcommand{\ie}{i.e.\ }
\newcommand{\cf}{cf.\ }
\newcommand{\resp}{respectively\ }
\newcommand{\Subsection}{\subsection}
\providecommand{\nobreakdash}{\nobreak\hbox{-}}
\setlength{\emergencystretch}{2em}
\multlinegap0pt
\usepackage{amssymb}
\usepackage{enumerate}
\usepackage{xcolor}
\newtheorem{corollary}{Corollary}
\newtheorem{lemma}{Lemma}
\newcommand{\f}{\mathcal{F}}
\title{Harmonicity of normal almost contact metric structures}
\author{M. Benyounes, T. Levasseur, E. Loubeau, E. Vergara-Diaz}
\date{Source: arXiv:2310.11089v2 (CC BY 4.0)}
\begin{document}
\maketitle

\noindent\emph{Source and licence.} Harmonicity of normal almost contact metric structures. Version arXiv:2310.11089v2 (CC BY 4.0), distributed by arXiv under the Creative Commons Attribution 4.0 International licence, http://creativecommons.org/licenses/by/4.0/, as recorded in the arXiv licence metadata for this exact version and shown on the abstract page https://arxiv.org/abs/2310.11089v2. Full licence notice: \texttt{licenses/arxiv-2310.11089-CC-BY-4.0.txt}.\par\smallskip

\noindent\textit{Editorial scope.} Counted unit: the article's lemma lemma3.2 (source0.tex lines 1127--1137). The article's complete original proof of it is delivered with it (source0.tex lines 1139--1212, one \texttt{proof} environment of 74 lines). No mathematics is written, repaired or re-derived: the statement and the proof are the article's own lines, in the article's own notation, and no retained line is substituted or re-flowed.  Retained uncounted as the source-local context the proof needs: lines 1070--1079; lines 1093--1098.  Nothing was omitted from any retained span, so the package carries no omission record.  No retained line contains a citation, so the wrapper carries no bibliography and no bibliography entry.  No retained line contains a figure environment, an image-inclusion command or drawing code, so no image is delivered and the package declares no asset dependency.  Editorial formatting only, in addition: an AMS \texttt{amsart} preamble that restates the article's own declarations and macros used by the retained lines, namely the environments \texttt{corollary}, \texttt{lemma} on separate counters, together with the standard packages those lines need.  The retained environments print in this document as Corollary 1, Lemma 1 in order of appearance, whereas the article prints the same items with the numbering of its own full document.  The complete native source of the article as distributed in the arXiv e-print for this exact version is bundled unchanged as \texttt{sources/148605/source0.tex}.
\begin{corollary}\label{corollary1}
  Let $(M^{2n+1},g)$ be a Riemannian manifold equipped with a
  normal almost contact metric structure
  $(\theta,\xi,\eta)$. Denote by $\f$ the contact sub-bundle,
  i.e.~the distribution orthogonal to $\xi$, by $\bar{J}$ the
  restriction of $\theta$ to $\f$ and by $\bar{\nabla}$ the
  connection induced on $\f$ by $\nabla$. Then
$$(\bar{\nabla}_{\bar{J}X}\bar{J})(\bar{J}Y) = (\bar{\nabla}_X \bar{J})(Y),$$
for any $X$ and $Y$ in $\f$.
\end{corollary}

\begin{lemma}\label{lemma2}
  Let $E$ be a horizontal vector field, that is a section of the
  horizontal sub-bundle $\f$, such that, at some point
  $x\in M^{2n+1}$, $(\bar{\nabla}E)(x)=0$. Then, at $x$,
$$[E,\bar{J}E]= (\bar{\nabla}_E \bar{J}) (E) +2 g(E, \nabla_{\theta E} \xi) \xi .$$
\end{lemma}

\begin{lemma}\label{lemma3.2}
  Let $E$ be a section of $\f$, such that, at some point
  $x\in M^{2n+1}$, $(\bar{\nabla}E)(x)=0$. Then, at $x$,
  $$
  [\bar{\nabla}_E \bar{\nabla}_E \bar{J},\bar{J}] = -2 [\bar{R}(E,\bar{J}E),\bar{J}
  ]-2\bar{\nabla}_{(\bar{\nabla}_E \bar{J})(E)}\bar{J}- [\bar{\nabla}_{\bar{J}E}
  \bar{\nabla}_{\bar{J}E} \bar{J},\bar{J}],
  $$
where $\bar{R}$ is the curvature tensor of the contact sub-bundle
$\f$ equipped with the connection $\bar{\nabla}$.
\end{lemma}

\begin{proof}
  Let $E$ be a section of $\f$ and $X$ a horizontal vector 
  extended locally such that $\bar{\nabla} X =0$ at the point
  $x\in M$ where we evaluate all expressions.  By the Leibniz
  rule, we have
  $$
  (\bar{\nabla}_E\bar{\nabla}_E \bar{J})(\bar{J}X)=\bar{\nabla}_E
  \big((\bar{\nabla}_E \bar{J})
  (\bar{J}X)\big)-(\bar{\nabla}_{E}\bar{J})((\bar{\nabla}_E \bar{J}) (X)).
  $$
  Using Corollary~\ref{corollary1} and Lemma~\ref{lemma2}, the
  first term on the right-hand side may be expressed in terms of
  the curvature tensor $\bar{R}$ as follows
\begin{align*}
&\bar{\nabla}_E\big((\bar{\nabla}_E \bar{J}) (\bar{J}X)\big)  = -\bar{\nabla}_E\big((\bar{\nabla}_{\bar{J}E}\bar{J}) (X)\big) 
 =-\bar{\nabla}_E\bar{\nabla}_{\bar{J}E}(\bar{J}X)+\bar{\nabla}_E(\bar{J}\bar{\nabla}_{\bar{J}E}
  X) \\
  &= -\bar{\nabla}_E\bar{\nabla}_{\bar{J}E}(\bar{J}X)+ \bar{J}(\bar{\nabla}_{E}\bar{\nabla}_{\bar{J}E}X)\\ 
& =-\bar{\nabla}_{\bar{J}E}\bar{\nabla}_E(\bar{J}X)+\bar{J}\bar{\nabla}_{\bar{J}E}\bar{\nabla}_E X -\bar{\nabla}_{[E,\bar{J}E]}(\bar{J}X)+ \bar{J}\bar{\nabla}_{[E,\bar{J}E]}X-[\bar{R}(E,\bar{J}E),\bar{J}](X)\\
&= -[\bar{R}(E,\bar{J}E),\bar{J}](X)  -\bar{\nabla}_{\bar{J}E}\bar{\nabla}_E(\bar{J}X)+\bar{J}\bar{\nabla}_{\bar{J}E}\bar{\nabla}_E  X - \bar{\nabla}_{(\bar{\nabla}_E \bar{J})(E)}(\bar{J}X)
  - 2g(  E, \nabla_{\bar{J} E}\xi) \bar{\nabla}_{\xi} (\bar{J}X)
  \\ & \phantom{xxx}+\bar{J}(\bar{\nabla}_{(\bar{\nabla}_E \bar{J})(E)}X)
+2g(  E, \nabla_{\bar{J} E}\xi) \bar{J}(\bar{\nabla}_{\xi}X)\\
  & =  -[\bar{R}(E,\bar{J}E),\bar{J}](X)  -\bar{\nabla}_{\bar{J}E}\bar{\nabla}_E(\bar{J}X)+\bar{J}\bar{\nabla}_{\bar{J}E}\bar{\nabla}_E X - \big(\bar{\nabla}_{(\bar{\nabla}_E \bar{J})(E)}\bar{J}\big)(X),
\end{align*}
since $\bar{\nabla}_{\xi} \bar{J} =0$. Now 
\begin{align*}
-\bar{\nabla}_{\bar{J}E}\bar{\nabla}_E(\bar{J}X) +\bar{J} (\bar{\nabla}_{\bar{J}E}\bar{\nabla}_E X)
&= -\bar{\nabla}_{\bar{J}E}((\bar{\nabla}_E \bar{J})(X)) -\bar{\nabla}_{\bar{J}E}(\bar{J}(\bar{\nabla}_E X))+\bar{J} (\bar{\nabla}_{\bar{J}E}\bar{\nabla}_E X)\\
&= -\bar{\nabla}_{\bar{J}E}((\bar{\nabla}_E \bar{J})(X))
= - \bar{\nabla}_{\bar{J}E}((\bar{\nabla}_{\bar{J}E}\bar{J})(\bar{J}X))\\
&= - (\bar{\nabla}_{\bar{J}E}\bar{\nabla}_{\bar{J}E}\bar{J})(\bar{J}X)- (\bar{\nabla}_{\bar{J}E}\bar{J}) (\bar{\nabla}_{\bar{J}E}(\bar{J}X))\\
&= - (\bar{\nabla}_{\bar{J}E}\bar{\nabla}_{\bar{J}E}\bar{J})(\bar{J}X)- (\bar{\nabla}_{\bar{J}E}\bar{J})((\bar{\nabla}_{\bar{J}E}\bar{J})(X)),
\end{align*}
because of the way we choose to extend the vector $X$.
Therefore
\begin{align}  
\bar{\nabla}_E \big((\bar{\nabla}_E \bar{J}) (\bar{J}X)\big) &=-[\bar{R}(E,\bar{J}E),\bar{J}](X)-\big(\bar{\nabla}_{(\bar{\nabla}_E\bar{J})(E)}\bar{J}\big)(X)\label{clubsuit}\\
&-(\bar{\nabla}_{\bar{J}E}\bar{\nabla}_{\bar{J}E}\bar{J}) (\bar{J}X) - (\bar{\nabla}_{\bar{J}E}\bar{J})\circ(\bar{\nabla}_{\bar{J}E}\bar{J})(X)\notag.
\end{align}
By definition of the covariant derivative of $\bar{\nabla}_E \bar{J}$
$$
\bar{\nabla}_E ((\bar{\nabla}_E \bar{J}) (\bar{J}X))= (\bar{\nabla}_E
\bar{\nabla}_E \bar{J})(\bar{J}X)+(\bar{\nabla}_{E}\bar{J})(
(\bar{\nabla}_{E}\bar{J})(X)),
$$
hence, using Equation~\eqref{clubsuit}, we have
\begin{align}
(\bar{\nabla}_E \bar{\nabla}_E \bar{J})(\bar{J}X) 
&=  -(\bar{\nabla}_{E}\bar{J})((\bar{\nabla}_{E}\bar{J})(X))+\bar{\nabla}_E ((\bar{\nabla}_E \bar{J}) (\bar{J}X))\notag\\
&=  -(\bar{\nabla}_{E}\bar{J})\circ(\bar{\nabla}_{E}\bar{J})(X)-[\bar{R}(E,\bar{J}E),\bar{J}](X)-(\bar{\nabla}_{(\bar{\nabla}_E\bar{J})(E)}\bar{J})(X)\notag\\
&-(\bar{\nabla}_{\bar{J}E}\bar{\nabla}_{\bar{J}E}\bar{J}) (\bar{J}X)  - (\bar{\nabla}_{\bar{J}E}\bar{J})\circ(\bar{\nabla}_{\bar{J}E}\bar{J})(X)  \notag \\
&= -2(\bar{\nabla}_{E}\bar{J})\circ(\bar{\nabla}_{E}\bar{J})(X)-[\bar{R}(E,\bar{J}E),\bar{J}](X)-(\bar{\nabla}_{(\bar{\nabla}_E\bar{J})(E)}\bar{J})(X)\label{eqstar}\\
&-(\bar{\nabla}_{\bar{J}E}\bar{\nabla}_{\bar{J}E}\bar{J}) (\bar{J}X) ,\notag
\end{align}
since, by Corollary~\ref{corollary1},
$$(\bar{\nabla}_{\bar{J}E}\bar{J})\circ(\bar{\nabla}_{\bar{J}E}\bar{J})(X) =
(\bar{\nabla}_{E}\bar{J})\circ(\bar{\nabla}_{E}\bar{J})(X).$$ 
% \begin{align*} 
% (\bar{\nabla}_{\bar{J}E}\bar{J})\circ(\bar{\nabla}_{\bar{J}E}\bar{J})(X)&= (\bar{\nabla}_{E}\bar{J})\circ(\bar{\nabla}_{E}\bar{J})(X).  
% \end{align*}
To compute the second term of $[\bar{\nabla}_E\bar{\nabla}_E \bar{J} ,
\bar{J}] (X)$ we proceed as follows: 
\begin{align*} 
(\bar{\nabla}_E \bar{\nabla}_E \bar{J})(\bar{J}^2X)&=-2\bar{J}\circ(\bar{\nabla}_{E}\bar{J})\circ(\bar{\nabla}_{E}\bar{J})(X)-[\bar{R}(E,\bar{J}E),\bar{J}](\bar{J}X)\\
&+\bar{J}(\bar{\nabla}_{(\bar{\nabla}_E\bar{J})(E)}\bar{J})(X)+(\bar{\nabla}_{\bar{J}E}\bar{\nabla}_{\bar{J}E}\bar{J}) (X),
\end{align*}
then
\begin{align*} 
-\bar{J}(\bar{\nabla}_E \bar{\nabla}_E \bar{J})(X)&=2(\bar{\nabla}_{E}\bar{J})\circ(\bar{\nabla}_{E}\bar{J})(X)-[\bar{R}(E,\bar{J}E),\bar{J}](X)\\
&-(\bar{\nabla}_{(\bar{\nabla}_E\bar{J})(E)}\bar{J})(X)+\bar{J}(\bar{\nabla}_{\bar{J}E}\bar{\nabla}_{\bar{J}E}\bar{J}) (X).
\end{align*}
Summing up this last equation with~\eqref{eqstar}, we obtain the result.
\end{proof}

\end{document}
